\documentclass[10pt,a4paper]{amsart}
\usepackage[margin=19mm]{geometry}
\usepackage{amsmath,amssymb,mathtools}
\usepackage[scr=rsfs,cal=euler]{mathalfa}
\usepackage{xfrac}
\usepackage[unicode,hidelinks,bookmarks=false]{hyperref}
\newcommand{\ie}{i.e.\ }
\newcommand{\cf}{cf.\ }
\newcommand{\resp}{respectively\ }
\newcommand{\Subsection}{\subsection}
\providecommand{\nobreakdash}{\nobreak\hbox{-}}
\setlength{\emergencystretch}{2em}
\multlinegap0pt
\usepackage{xcolor}
\usepackage{amssymb}
\usepackage[shortlabels]{enumitem}
\usepackage{float}
\newtheorem{lem}{}
\usepackage{cleveref}
\crefname{lem}{}{s}
\newcommand{\N}{\mathbb{N}}
\newcommand{\Z}{\mathbb{Z}}
\title{Sets of large values of polynomial multi-correlation functions}
\author{Vitaly Bergelson, Rigoberto Zelada}
\date{Source: arXiv:2605.23050v1 (CC BY 4.0)}
\begin{document}
\maketitle

\noindent\emph{Source and licence.} Sets of large values of polynomial multi-correlation functions. Version arXiv:2605.23050v1 (CC BY 4.0), distributed by arXiv under the Creative Commons Attribution 4.0 International licence, http://creativecommons.org/licenses/by/4.0/, as recorded in the arXiv licence metadata for this exact version and shown on the abstract page https://arxiv.org/abs/2605.23050v1. Full licence notice: \texttt{licenses/arxiv-2605.23050-CC-BY-4.0.txt}.\par\smallskip

\noindent\textit{Editorial scope.} Counted unit: the article's lem 6.lem:DegreeOfLinearRestriction (source0.tex lines 1098--1102). The article's complete original proof of it is delivered with it (source0.tex lines 1103--1116, one \texttt{proof} environment of 14 lines). No mathematics is written, repaired or re-derived: the statement and the proof are the article's own lines, in the article's own notation, and no retained line is substituted or re-flowed.  No further source-local context is needed: every cross-reference of every retained line resolves inside the two delivered spans.  Nothing was omitted from any retained span, so the package carries no omission record.  No retained line contains a citation, so the wrapper carries no bibliography and no bibliography entry.  No retained line contains a figure environment, an image-inclusion command or drawing code, so no image is delivered and the package declares no asset dependency.  Editorial formatting only, in addition: an AMS \texttt{amsart} preamble that restates the article's own declarations and macros used by the retained lines, namely the environments \texttt{lem} on separate counters, together with the standard packages those lines need.  The retained environments print in this document as Lemma 1 in order of appearance, whereas the article prints the same items with the numbering of its own full document.  The complete native source of the article as distributed in the arXiv e-print for this exact version is bundled unchanged as \texttt{sources/201178/source0.tex}.
\begin{lem}\label{6.lem:DegreeOfLinearRestriction}
    Let $d\in\N$ and let $p\in \Z[x_1,...,x_d]$ be a non-constant polynomial. There is a set $E\subseteq \Z^d$ with uniform density one such that for any $\vec a=(a_1,...,a_d)\in E$ the polynomial 
    $$q_{\vec a}(x)=p(a_1x,...,a_dx)$$
    satisfies $\deg_x(q_{\vec a})=\deg_{x_1,...,x_d}(p)$.
\end{lem}

\begin{proof}[Proof of \cref{6.lem:DegreeOfLinearRestriction}] Let $D=\deg_{x_1,...,x_d}(p)$ and   write
$$
p(x_1,...,x_d)=\sum_{s=0}^D\sum_{\substack{j_1+\cdots+j_d=s\\j_1,...,j_d\in\N\cup\{0\}}}c_{j_1,...,j_d}x_1^{j_1}\cdots x_d^{j_d}.
$$
Consider the polynomial 
$$
u(x_1,...,x_d):= \sum_{\substack{j_1+\cdots+j_d=D\\j_1,...,j_d\in\N\cup\{0\}}}c_{j_1,...,j_d}x_1^{j_1}\cdots x_d^{j_d}
$$
and note that, like $p$, $u$ is a non-constant polynomial. Thus,  the set $F:=\{\vec a\in\Z^d\,|\,u(\vec a)=0 \}$ satisfies $d^*(F)=0$. Letting $E=\Z^d\setminus F$ it now follows that for every $\vec a=(a_1,...,a_d)\in E$, the polynomial  
$$
q_{\vec a}(x)=\sum_{s=0}^D\left(\sum_{\substack{j_1+\cdots+j_d=s\\j_1,...,j_d\in\N\cup\{0\}}}c_{j_1,...,j_d}a_1^{j_1}\cdots a_d^{j_d}\right)x^s.
$$
satisfies $\deg_x(q_{\vec a})=D$.
\end{proof}

\end{document}
